\documentclass[11pt,a4paper]{article}

\usepackage[T1]{fontenc}
\usepackage[utf8]{inputenc}
\usepackage[a4paper,margin=2.5cm]{geometry}
\usepackage[hidelinks]{hyperref}
\usepackage{enumitem}
\usepackage{longtable}

\setlist{noitemsep,topsep=4pt}

\title{Monad Fleet Device Protocol Specification\\fleet.v2 Implementation Profile}
\author{FleetManager Project}
\date{\today}

\begin{document}
\maketitle

\section{Document Purpose}
This document is a full, implementation-first specification of the currently deployed Monad Fleet device protocol profile.
It is intended for engineers building new device agents (Raspberry Pi or other Wi-Fi-capable devices) that must join the system and execute experiments safely.

The scope includes:
\begin{itemize}
  \item gRPC contract (\texttt{fleet.v2}),
  \item run lifecycle and idempotency,
  \item device-side persistence model,
  \item metrics forwarding pipeline (Prometheus/Mimir/Grafana),
  \item eLabFTW integration profile used by the running server.
\end{itemize}

\textbf{Implementation baseline:} February 18, 2026.
This document describes the behavior currently implemented in this repository, not a hypothetical future API.

\section{Normative Language}
The words \textbf{MUST}, \textbf{SHOULD}, and \textbf{MAY} are used in the RFC-style sense:
\begin{itemize}
  \item \textbf{MUST}: required for interoperability with this profile.
  \item \textbf{SHOULD}: strongly recommended; valid reasons may exist to deviate.
  \item \textbf{MAY}: optional.
\end{itemize}

\section{Source of Truth}
\begin{itemize}
  \item v2 protobuf: \texttt{monad-fleet-service/proto/fleet\_gateway\_v2.proto}
  \item server implementation: \texttt{monad-fleet-service/gateway\_server.py}
  \item reference client: \texttt{device-sim/agent\_v2\_client.py}
  \item legacy v1 protobuf (still served): \texttt{monad-fleet-service/proto/fleet\_gateway.proto}
\end{itemize}

\section{System Roles and Boundaries}
\subsection{Roles}
\begin{itemize}
  \item \textbf{Agent}: software on measurement device; executes policy and reports run results.
  \item \textbf{Monad Fleet}: orchestration service; resolves assignments, serves policy, ingests reports.
  \item \textbf{eLabFTW}: experiment registry and scheduler/event storage backend.
  \item \textbf{Metrics stack (optional)}: Prometheus + Mimir + Grafana.
\end{itemize}

\subsection{Boundary contract}
Agents communicate with Monad Fleet only.
Agents do \textbf{not} need eLabFTW-specific knowledge.
Backend routing (eLabFTW, time-series, object storage) is a server deployment concern.

\section{Architecture Overview}
\begin{verbatim}
Device Agent <-> gRPC fleet.v2 <-> Monad Fleet <-> eLabFTW API
          \-> optional HTTP ingest (/ingest/v1/metrics, /ingest/v1/artifacts)
                                           -> /metrics -> Prometheus -> Mimir -> Grafana
\end{verbatim}

\section{Transport and Endpoints}
\begin{itemize}
  \item gRPC transport: HTTP/2
  \item gRPC default port: \texttt{50060}
  \item reference client transport today: \texttt{grpc.insecure\_channel()} (no TLS)
  \item HTTP sidecar port: \texttt{9108}
  \item HTTP endpoints:
    \begin{itemize}
      \item \texttt{GET /healthz}
      \item \texttt{GET /metrics}
      \item \texttt{POST /ingest/v1/metrics}
      \item \texttt{POST /ingest/v1/artifacts}
    \end{itemize}
\end{itemize}

\textbf{Security note:} run on trusted LAN/VPN for this profile, or add mTLS in a hardened profile.

\section{Protocol Design Principles}
\subsection{Backend agnostic send operation}
The agent send operation is \texttt{PublishReport} (plus optional \texttt{PublishEvents}).
Reports carry:
\begin{itemize}
  \item structured events,
  \item artifact references,
  \item summary metrics.
\end{itemize}

This keeps the device contract generic while allowing server-side forwarding to different destinations.

\subsection{Forward compatibility}
\begin{itemize}
  \item Agents \textbf{MUST} ignore unknown fields and enum values.
  \item Servers \textbf{MUST} tolerate omitted optional fields.
  \item \texttt{Event.metrics} and \texttt{Report.summary\_metrics} are open extension maps.
\end{itemize}

\section{Service Definition}
Served package: \texttt{fleet.v2}. Service: \texttt{FleetManager}.

\begin{verbatim}
rpc Hello(HelloRequest) returns (HelloResponse);
rpc GetAssignment(GetAssignmentRequest) returns (GetAssignmentResponse);
rpc GetPolicy(GetPolicyRequest) returns (GetPolicyResponse);
rpc AckPrepared(AckPreparedRequest) returns (AckPreparedResponse);
rpc PublishReport(PublishReportRequest) returns (PublishReportResponse);
rpc PublishEvents(PublishEventsRequest) returns (PublishEventsResponse); // optional
\end{verbatim}

\section{Run Lifecycle}
\subsection{Operational phases}
\begin{itemize}
  \item \textbf{PREPARE}: \texttt{Hello}, \texttt{GetAssignment}, \texttt{GetPolicy}, \texttt{AckPrepared}
  \item \textbf{MEASURE}: local policy command execution on agent
  \item \textbf{REPORT}: \texttt{PublishReport}; optional \texttt{PublishEvents}
\end{itemize}

\subsection{Recommended loop}
\begin{verbatim}
1) Hello(agent descriptor)
2) repeat:
   a) GetAssignment(agent_id)
   b) if NO_WORK: sleep(recommended_prepare_poll_sec)
   c) GetPolicy(agent_id, experiment_id, last_policy_id)
   d) if NOT_MODIFIED: sleep(recommended_prepare_poll_sec)
   e) AckPrepared(...)
   f) execute command_groups
   g) PublishReport(run_id, events, artifacts, summary)
   h) on ACCEPTED or DUPLICATE: cache policy_id and continue
\end{verbatim}

\section{Control Plane Mode Semantics}
\subsection{Modes}
\begin{itemize}
  \item \texttt{RF\_SHARING}: expected for Wi-Fi-only devices.
  \item \texttt{DUAL\_NIC}: allowed only when Ethernet path is available and validated.
\end{itemize}

\subsection{Current server behavior}
\begin{itemize}
  \item Requesting \texttt{DUAL\_NIC} without Ethernet advertisement is downgraded to \texttt{RF\_SHARING}.
  \item \texttt{AckPrepared} rejects \texttt{DUAL\_NIC} when:
    \begin{itemize}
      \item \texttt{route\_verified=false}, or
      \item \texttt{control\_plane\_iface} is provided but is not Ethernet-like (\texttt{eth*}).
    \end{itemize}
\end{itemize}

\section{RPC Semantics}
\subsection{\texttt{Hello}}
\textbf{Required:} \texttt{agent.agent\_id}

\textbf{Purpose:}
\begin{itemize}
  \item register/refresh device presence,
  \item update descriptor metadata,
  \item negotiate effective control-plane mode.
\end{itemize}

\textbf{Important response fields:}
\begin{itemize}
  \item \texttt{allowed\_mode}, \texttt{mode\_reason}
  \item \texttt{recommended\_prepare\_poll\_sec}
  \item \texttt{server\_version}, \texttt{server\_time}, \texttt{max\_batch\_size}
\end{itemize}

\subsection{\texttt{GetAssignment}}
\textbf{Required:} \texttt{agent\_id}

\textbf{Responses:}
\begin{itemize}
  \item \texttt{ASSIGNED}: returns \texttt{experiment\_id}, \texttt{policy\_id}, measurement window.
  \item \texttt{NO\_WORK}: no currently selected policy for this agent.
\end{itemize}

\subsection{\texttt{GetPolicy}}
\textbf{Required:} \texttt{agent\_id}. Optional: \texttt{experiment\_id}, \texttt{last\_policy\_id}.

\textbf{Responses:}
\begin{itemize}
  \item \texttt{OK}: full \texttt{Policy}.
  \item \texttt{NOT\_MODIFIED}: \texttt{last\_policy\_id} still current.
  \item \texttt{NO\_POLICY}: no policy found or mismatch.
\end{itemize}

\subsection{\texttt{AckPrepared}}
\textbf{Required:} \texttt{agent\_id}, \texttt{preparation\_id}, \texttt{policy\_id}

\textbf{Validation:}
\begin{itemize}
  \item reject if no assigned policy,
  \item reject on policy revision mismatch,
  \item reject invalid \texttt{DUAL\_NIC} checks.
\end{itemize}

\textbf{Idempotency:}
duplicate \texttt{preparation\_id} is accepted and treated as already acknowledged.

\subsection{\texttt{PublishReport}}
\textbf{Required:}
\begin{itemize}
  \item \texttt{report.run\_id}
  \item agent identity from \texttt{agent\_id} or \texttt{report.agent\_id}
\end{itemize}

\textbf{Current server behavior:}
\begin{itemize}
  \item map v2 report/events into internal event model,
  \item write run state to eLabFTW scheduler event metadata,
  \item process artifact references as \texttt{ARTIFACT\_UPLOADED} events,
  \item export numeric summary/event metrics to \texttt{monad\_fleet\_metric\_value}.
\end{itemize}

\textbf{Status values:}
\begin{itemize}
  \item \texttt{ACCEPTED}
  \item \texttt{DUPLICATE} (same \texttt{run\_id|agent\_id|policy\_id})
  \item \texttt{REJECTED}
\end{itemize}

\subsection{\texttt{PublishEvents}}
In this deployment profile, live events are disabled (\texttt{ALLOW\_LIVE\_EVENTS=false}).
When disabled, each ack is \texttt{REJECTED} with reason \texttt{live events disabled}.

\section{Idempotency and De-duplication}
\begin{itemize}
  \item prepare dedupe key: \texttt{preparation\_id}
  \item event dedupe key: \texttt{event\_id}
  \item report dedupe key: \texttt{run\_id|agent\_id|policy\_id}
\end{itemize}

Server runtime state is persisted at:
\begin{itemize}
  \item \texttt{/data/state.json} (container path)
\end{itemize}

\section{Data Model Reference}
\subsection{Common conventions}
\begin{itemize}
  \item Time on wire: \texttt{google.protobuf.Timestamp} in UTC.
  \item Policy JSON timestamps: ISO-8601 UTC (\texttt{Z} suffix).
  \item IDs are opaque strings; recommended:
    \begin{itemize}
      \item \texttt{agent\_id}: stable MAC or UUID
      \item \texttt{run\_id}: UUIDv4
      \item \texttt{event\_id}: UUIDv4 or deterministic hash
      \item \texttt{policy\_id}: server-generated immutable revision (\texttt{sha256:<digest>})
    \end{itemize}
  \item Numeric samples in maps are encoded as strings and parsed to float server-side.
\end{itemize}

\subsection{\texttt{AgentDescriptor}}
\begin{longtable}{p{4.1cm}p{1.6cm}p{9.3cm}}
\textbf{Field} & \textbf{Req} & \textbf{Meaning} \\
\hline
\endhead
\texttt{agent\_id} & yes & Stable unique device identifier. \\
\texttt{hostname} & no & Human-readable host name. \\
\texttt{agent\_version} & no & Agent software version. \\
\texttt{interfaces} & no & Advertised interfaces for audit/mode checks. \\
\texttt{requested\_mode} & no & Requested control-plane mode. \\
\texttt{control\_plane\_iface} & no & Interface intended for control plane. \\
\texttt{fleet\_manager\_target} & no & Host:port used for route verification context. \\
\texttt{capabilities} & no & Free-form capability tags (\texttt{wifi}, \texttt{ble}, \texttt{csi}, ...). \\
\hline
\end{longtable}

\subsection{\texttt{Policy}}
\begin{longtable}{p{4.1cm}p{1.6cm}p{9.3cm}}
\textbf{Field} & \textbf{Req} & \textbf{Meaning} \\
\hline
\endhead
\texttt{experiment\_id} & yes & Experiment identity (\texttt{elabftw:<id>} in current deployment). \\
\texttt{policy\_id} & yes & Immutable revision id for caching/idempotency. \\
\texttt{issued\_at} & no & Issuance time. \\
\texttt{measurement\_from/to} & no & Intended run window. \\
\texttt{command\_groups} & yes & Ordered groups of commands. \\
\texttt{allowed\_mode} & no & Effective authorized mode. \\
\texttt{control\_plane\_iface} & no & Effective iface for \texttt{DUAL\_NIC}. \\
\hline
\end{longtable}

\subsection{\texttt{CommandGroup} and \texttt{Command}}
\begin{longtable}{p{4.1cm}p{1.6cm}p{9.3cm}}
\textbf{Field} & \textbf{Req} & \textbf{Meaning} \\
\hline
\endhead
\texttt{CommandGroup.id} & yes & Group correlation id. \\
\texttt{CommandGroup.name} & no & Display name. \\
\texttt{CommandGroup.from/to} & no & Optional per-group window. \\
\texttt{CommandGroup.failure\_mode} & no & \texttt{FAIL\_FAST} or \texttt{CONTINUE\_ON\_ERROR}. \\
\texttt{CommandGroup.commands} & yes & Ordered command list. \\
\texttt{Command.id} & yes & Command correlation id. \\
\texttt{Command.type} & yes & \texttt{SHELL}, \texttt{WIFI\_SCAN}, \texttt{BLE\_SCAN}, \texttt{CAPTURE\_CSI}. \\
\texttt{Command.cmdline} & no & Command line (primarily for \texttt{SHELL}). \\
\texttt{Command.argv} & no & Optional argv form. \\
\texttt{Command.env} & no & Environment key-values. \\
\texttt{Command.timeout\_ms} & no & Command timeout in ms. \\
\texttt{Command.retry} & no & Retry policy. \\
\texttt{Command.expected\_artifacts} & no & Expected outputs. \\
\hline
\end{longtable}

\subsection{\texttt{RetryPolicy}}
\begin{longtable}{p{4.1cm}p{1.6cm}p{9.3cm}}
\textbf{Field} & \textbf{Req} & \textbf{Meaning} \\
\hline
\endhead
\texttt{max\_attempts} & no & Total attempts including first attempt. \\
\texttt{backoff\_ms} & no & Fixed wait between attempts. \\
\hline
\end{longtable}

\subsection{\texttt{Report} and \texttt{Event}}
\begin{longtable}{p{4.1cm}p{1.6cm}p{9.3cm}}
\textbf{Field} & \textbf{Req} & \textbf{Meaning} \\
\hline
\endhead
\texttt{Report.run\_id} & yes & Unique run identifier. \\
\texttt{Report.experiment\_id} & yes & Experiment id of this run. \\
\texttt{Report.policy\_id} & yes & Policy revision id of this run. \\
\texttt{Report.agent\_id} & no & Optional when request-level \texttt{agent\_id} is provided. \\
\texttt{Report.started\_at/finished\_at} & no & Informational run timing. \\
\texttt{Report.status} & yes & Human-readable final status (\texttt{OK}/\texttt{FAILED}/...). \\
\texttt{Report.events} & no & Structured event list. \\
\texttt{Report.artifacts} & no & Artifact references only. \\
\texttt{Report.summary\_metrics} & no & Free map; numeric values exported to metrics backend. \\
\texttt{Event.event\_id} & no & Event unique id (server can synthesize). \\
\texttt{Event.type} & yes & Run/command/artifact/error type enum. \\
\texttt{Event.timestamp} & no & Event timestamp. \\
\texttt{Event.group\_id/command\_id} & no & Correlation fields. \\
\texttt{Event.exit\_code} & no & Exit code for command events. \\
\texttt{Event.duration\_ms} & no & Duration for command events. \\
\texttt{Event.message} & no & Short text payload. \\
\texttt{Event.metrics} & no & Free map; numeric values exported to metrics backend. \\
\hline
\end{longtable}

\subsection{\texttt{ArtifactRef}}
\begin{longtable}{p{4.1cm}p{1.6cm}p{9.3cm}}
\textbf{Field} & \textbf{Req} & \textbf{Meaning} \\
\hline
\endhead
\texttt{name} & yes & Artifact label. \\
\texttt{uri} & yes & Artifact location reference (\texttt{spool://}, \texttt{s3://}, \texttt{https://}...). \\
\texttt{size\_bytes} & no & Artifact size for audit. \\
\texttt{sha256} & no & Artifact checksum. \\
\hline
\end{longtable}

\section{Policy Profile Used in This Deployment}
\subsection{Policy source}
Fleet policies are loaded from eLabFTW experiments tagged \texttt{fleet}.
Categories are not used by the fleet selector.

\textbf{Required authoring mode:}
store policy in experiment metadata JSON editor (not as attachment).

\textbf{Metadata keys searched in order:}
\begin{verbatim}
fleet.policy, policy, fleet_policy, policy_json
\end{verbatim}

\subsection{Selection logic}
\begin{itemize}
  \item candidate set: experiments with \texttt{fleet} tag
  \item optional tag matching by \texttt{RESOURCE\_TAG\_PREFIX} (default \texttt{device:})
  \item optional \texttt{target\_selector} on \texttt{device\_ids/device\_types/locations}
  \item newest modified valid policy wins
\end{itemize}

\subsection{Policy revision id}
Server canonicalizes policy JSON and computes:
\begin{verbatim}
policy_id = sha256:<digest>
\end{verbatim}

\subsection{Reference policy example}
\begin{verbatim}
{
  "fleet": {
    "policy": {
      "target_selector": {
        "device_ids": ["2c:cf:67:80:f5:d9"]
      },
      "range": {
        "from": "2026-02-17T19:41:31Z",
        "to":   "2026-02-17T21:42:31Z"
      },
      "command_groups": [
        {
          "id": "wifi-ble-csi-v3",
          "name": "WiFi BLE CSI v3",
          "failure_mode": "FAIL_FAST",
          "commands": [
            {"id": "wifi-scan",  "type": "WIFI_SCAN",   "timeout_s": 15},
            {"id": "ble-scan",   "type": "BLE_SCAN",    "timeout_s": 15},
            {"id": "csi-capture","type": "CAPTURE_CSI", "timeout_s": 15}
          ]
        }
      ]
    }
  }
}
\end{verbatim}

\subsection{RSSI sampling policy example (30s, 1s interval)}
\begin{verbatim}
{
  "fleet": {
    "policy": {
      "target_selector": {"device_ids": ["2c:cf:67:80:f5:d9"]},
      "command_groups": [
        {
          "id": "wifi-rssi-series",
          "name": "WiFi RSSI Sampling 30s/1s",
          "failure_mode": "FAIL_FAST",
          "commands": [
            {
              "id": "wifi-rssi-series",
              "type": "WIFI_SCAN",
              "timeout_s": 35,
              "env": {
                "WIFI_SAMPLE_DURATION_S": "30",
                "WIFI_SAMPLE_INTERVAL_S": "1",
                "WIFI_SAMPLE_EMIT_HTTP": "1"
              }
            }
          ]
        }
      ]
    }
  }
}
\end{verbatim}

This profile lets experiment authors define sampling cadence in the policy itself.
Device-side constraints may clamp requested values (minimum interval, maximum duration, maximum points).

\section{Device Persistence and Deferred Upload}
\subsection{Goal}
Preserve run outputs during network outages and retry upload later without data loss.

\subsection{Local spool layout}
\begin{verbatim}
DATA_ROOT/
  pending/<run_id>/
    context.json
    events.ndjson
    report.json
    artifacts/
      <command-id>-<timestamp>.log
      run-summary.json
  sent/<run_id>/
    ...same structure...
\end{verbatim}

\subsection{Upload semantics}
\begin{itemize}
  \item agent writes completed run bundle under \texttt{pending/<run\_id>/}
  \item agent retries \texttt{PublishReport} for pending runs
  \item on \texttt{ACCEPTED} or \texttt{DUPLICATE}, run is moved to \texttt{sent/<run\_id>/}
  \item sent history is pruned by \texttt{SENT\_RETENTION\_DAYS} (default 14)
\end{itemize}

\section{Artifact Handling Profile}
Artifacts are persisted locally first, then published through two layers:
\begin{itemize}
  \item \texttt{PublishReport} carries \texttt{ArtifactRef} metadata,
  \item optional HTTP upload (\texttt{POST /ingest/v1/artifacts}) sends local files to eLabFTW experiment uploads.
\end{itemize}

When HTTP artifact upload succeeds, agent rewrites artifact URI from local spool form to:
\begin{verbatim}
elabftw://experiments/<id>/uploads/<upload_id>
\end{verbatim}

If upload fails or is disabled, URI remains \texttt{spool://...} and run reporting still succeeds.

Relevant artifact classes for experiment analysis:
\begin{itemize}
  \item Wi-Fi: scan/link snapshots, command logs
  \item BLE: scan output snapshots, command logs
  \item CSI: capture logs, collector outputs, or reference notes
  \item run-level summary (\texttt{run-summary.json})
\end{itemize}

Current size guardrail:
\begin{itemize}
  \item server-side endpoint limit: \texttt{ARTIFACT\_MAX\_BYTES} (default 20 MiB)
  \item agent-side upload limit: \texttt{ARTIFACT\_UPLOAD\_MAX\_BYTES} (default 20 MiB)
\end{itemize}

\section{Metrics Pipeline Profile (v3)}
\subsection{Objective}
Store numeric Wi-Fi/BLE/CSI-derived metrics in time-series backend (Mimir) and visualize in Grafana.

\subsection{Path}
\begin{verbatim}
Agent summary/event metrics -> PublishReport -> Fleet /metrics -> Prometheus
                                                               -> remote_write -> Mimir -> Grafana
Low-resource sender JSON   -> /ingest/v1/metrics -> Fleet /metrics -> ...
\end{verbatim}

\subsection{Low-resource sender payload}
\begin{verbatim}
POST /ingest/v1/metrics
{
  "source": "embedded-board",
  "device_id": "lowlevel-board-01",
  "metrics": [
    {"name": "wifi_rssi_dbm", "value": -58.2},
    {"name": "ble_adv_count", "value": 81},
    {"name": "csi_frames_count", "value": 420}
  ]
}
\end{verbatim}
\textbf{Note:} this payload is optional for development/demo paths. Real-only smoke runs can disable synthetic injection entirely.

\subsection{Metrics vs raw captures}
Mimir stores numeric samples, not binary capture blobs.
Therefore:
\begin{itemize}
  \item raw CSI/BLE/PCAP data should be persisted as artifacts (local/object storage and optionally mirrored to eLabFTW uploads)
  \item derived numeric metrics should be sent via event/summary maps
\end{itemize}
\textbf{Current deployment profile:} CSI numeric keys (\texttt{csi\_*}) are excluded from time-series export to Prometheus/Mimir/Grafana.
CSI remains artifact-first (run artifacts and optional eLabFTW uploads).

\section{eLabFTW Mapping Profile}
\subsection{Experiment authoring}
\begin{enumerate}
  \item create experiment
  \item add tag \texttt{fleet}
  \item store policy in metadata JSON (\texttt{fleet.policy})
  \item optionally use \texttt{device:<tag>} matching tags
\end{enumerate}

\subsection{Device resources}
On first \texttt{Hello}, server ensures resource item exists for \texttt{Device <agent\_id>} and updates metadata.
Bookable defaults are enforced in current deployment:
\begin{itemize}
  \item \texttt{is\_bookable=1}
  \item \texttt{book\_max\_minutes=180}
  \item \texttt{book\_can\_overlap=1}
  \item \texttt{book\_is\_cancellable=1}
\end{itemize}

\subsection{Scheduler run mapping}
Run aggregation key:
\begin{verbatim}
experiment_id | policy_revision | device_id
\end{verbatim}

Stored metadata includes status, event history, summary metrics, and last event payload.
Terminal status is treated as monotonic and should not be downgraded by later artifact events.

\section{Raspberry Pi Reference Profile}
\subsection{Validated hardware reference}
Current validated setup used in smoke runs:
\begin{itemize}
  \item \texttt{Raspberry Pi 5 Model B Rev 1.0}
  \item SoC family: \texttt{Broadcom BCM2712}
  \item PCIe M.2 Wi-Fi card: \texttt{Intel Wi-Fi 6E AX210/AX1675 2x2} (\texttt{8086:2725})
\end{itemize}

\subsection{Reference scripts}
\begin{verbatim}
scripts/rpi/setup_ssh_key.sh
scripts/rpi/deploy_agent.sh
scripts/rpi/install_systemd_service.sh
scripts/rpi/smoke_test_agent.sh
\end{verbatim}

\subsection{Recommended environment}
\begin{verbatim}
FLEET_MANAGER_HOST=<server_host>
FLEET_MANAGER_PORT=50060
AGENT_ID=<stable_mac_or_uuid>
CONTROL_PLANE_MODE=RF_SHARING
CONTROL_PLANE_IFACE=wlan0
WIFI_SCAN_IFACE=wlan0
EXECUTE_POLICY=true
CAPABILITIES=csi,ble,wifi
DATA_ROOT=./data
SENT_RETENTION_DAYS=14

# Optional device-side sampling constraints (applied to WIFI_SCAN sampling mode)
WIFI_SAMPLE_MIN_INTERVAL_S=1
WIFI_SAMPLE_MAX_DURATION_S=300
WIFI_SAMPLE_MAX_POINTS=600
ENABLE_HTTP_SAMPLE_INGEST=true

# Extended real-device observability
ENABLE_WIFI_EXTENDED_METRICS=true
ENABLE_DEVICE_STATUS_METRICS=true

# Optional CSI collector hook (recommended for real CSI)
CSI_COLLECTOR_CMD="<collector command line>"
CSI_OUTPUT_PATH="<single output file path>"
CSI_OUTPUT_GLOB="<glob(s) for collector outputs>"
CSI_FRAMES_REGEX="<regex with one numeric capture>"
CSI_OUTPUT_MAX_FILES=5
CSI_PARSE_MAX_BYTES=1048576
DISABLE_CSI_CAPTURE=false
\end{verbatim}

\textbf{Note:} real Wi-Fi/BLE metrics depend on local tools (\texttt{iw}, \texttt{bluetoothctl}) and radio access.
\textbf{CSI note:} this interface now supports an external CSI collector hook. Without configured collector command/output, CSI metrics will remain at zero and report explicit collector-status metrics.
\textbf{Control-plane safety note:} when the current route to Fleet Manager is over Wi-Fi, active CSI collection can disrupt connectivity. The smoke tooling defaults to \texttt{DISABLE\_CSI\_CAPTURE=true} in this case unless explicitly overridden.

\subsection{CSI collector profile}
For \texttt{CAPTURE\_CSI}, the agent now supports two modes:
\begin{itemize}
  \item direct command from policy (\texttt{cmdline}/\texttt{argv})
  \item environment-driven collector hook (\texttt{CSI\_COLLECTOR\_CMD}, optional output file/glob settings)
\end{itemize}

Returned CSI status metrics include:
\begin{itemize}
  \item \texttt{csi\_collector\_configured}
  \item \texttt{csi\_collector\_exit\_code}
  \item \texttt{csi\_capture\_ok}
  \item \texttt{csi\_capture\_disabled}
  \item \texttt{csi\_frames\_count}
  \item \texttt{csi\_output\_files\_count}
  \item \texttt{csi\_output\_bytes\_total}
\end{itemize}

Collector output files can be mirrored as run artifacts and uploaded to eLabFTW (subject to size limits).
\textbf{Hardware caveat:} Intel AX210 with stock Linux drivers does not expose CSI frames by default. Real CSI requires a compatible collector/driver stack; until configured, \texttt{csi\_collector\_configured=0} and \texttt{csi\_frames\_total=0}.

\section{Smoke Test From Clean Fleet/Metrics State}
\subsection{Automated script}
\begin{verbatim}
scripts/smoke/v3_end_to_end_smoke.sh
\end{verbatim}

\subsection{Manual reset helper}
\begin{verbatim}
scripts/reset/dev_reset.sh
\end{verbatim}

\subsection{Reset behavior}
The smoke flow is designed to preserve existing user data in eLabFTW/MySQL.
It resets fleet runtime and metrics storage state only.

\textbf{Does reset:}
\begin{itemize}
  \item \texttt{./data/fleet/state.json}
  \item \texttt{./data/fleet/ingest-metrics.ndjson}
  \item \texttt{./data/fleet/ingest-artifacts.ndjson}
  \item Prometheus local TSDB volume
  \item Mimir local TSDB volume
\end{itemize}

\textbf{Does not reset:}
\begin{itemize}
  \item eLabFTW user accounts
  \item existing MySQL experiment history
\end{itemize}

\subsection{Smoke parameters}
\begin{verbatim}
DEVICE_MODEL=02:42:ac:14:00:04
DEVICE_PI=2c:cf:67:80:f5:d9
RESET_STATE=true|false
RESET_METRICS=true|false
RESET_GRAFANA=true|false
RUN_PI=true|false
PI_HOST=monad-rpi5.local
DO_BUILD=true|false
CLEANUP=true|false

# Optional smoke profile
SMOKE_PROFILE=full|rssi|wireless_spec
RSSI_DURATION_S=30
RSSI_INTERVAL_S=1
RSSI_INTERVAL_S_PI=5
RSSI_ARTIFACT_STRIDE=0

# Optional 2-minute full profile (Wi-Fi sampling + BLE + CSI artifacts)
FULL_DURATION_S=120
FULL_INTERVAL_S=2
FULL_INTERVAL_S_PI=5
FULL_ARTIFACT_STRIDE=0
FULL_BLE_TIMEOUT_S=10
FULL_CSI_TIMEOUT_S=10
REPORT_RPC_TIMEOUT_S=120
CONTROL_RPC_TIMEOUT_S=30
HELLO_RPC_TIMEOUT_S=<optional; defaults to CONTROL_RPC_TIMEOUT_S>
ASSIGNMENT_RPC_TIMEOUT_S=<optional; defaults to CONTROL_RPC_TIMEOUT_S>
POLICY_RPC_TIMEOUT_S=<optional; defaults to CONTROL_RPC_TIMEOUT_S>
ACK_PREPARED_RPC_TIMEOUT_S=<optional; defaults to CONTROL_RPC_TIMEOUT_S>

# Wireless sensing spec profile (short, realistic Pi run)
WIRELESS_SPEC_DURATION_S=30
WIRELESS_SPEC_INTERVAL_S=2
WIRELESS_SPEC_INTERVAL_S_PI=3
WIRELESS_SPEC_ARTIFACT_STRIDE=0
WIRELESS_SPEC_BLE_TIMEOUT_S=12
WIRELESS_SPEC_CSI_TIMEOUT_S=15
WIRELESS_SPEC_ENV_SNAPSHOT=true

# Optional CSI collector pass-through (for real AX210 CSI if configured)
CSI_COLLECTOR_CMD="<collector command>"
CSI_OUTPUT_PATH="<single output file>"
CSI_OUTPUT_GLOB="<glob for output files>"
CSI_FRAMES_REGEX="<regex with one numeric capture group>"
CSI_OUTPUT_MAX_FILES=5
CSI_PARSE_MAX_BYTES=1048576
ALLOW_WIFI_DISRUPTIVE_CSI=true|false
DISABLE_CSI_CAPTURE=true|false

# Optional artifact upload controls
ENABLE_ELAB_ARTIFACT_UPLOAD=true|false
ARTIFACT_UPLOAD_MAX_BYTES=20971520

# Optional demo-only synthetic metric injection
INJECT_HYPOTHETICAL_METRICS=true|false
\end{verbatim}

For Raspberry Pi runs, a higher interval is recommended (for example 5s or 10s), depending on radio/tooling overhead.
\textbf{Real-only recommendation:} run with \texttt{RUN\_PI=true} and \texttt{INJECT\_HYPOTHETICAL\_METRICS=false}.
\textbf{Upload stability note:} sampled Wi-Fi commands persist first+last sample artifacts by default (\texttt{*\_ARTIFACT\_STRIDE=0}) to keep report size bounded; use stride \texttt{1} if full per-sample artifacts are explicitly needed.
\textbf{Automation note:} Pi smoke scripts use SSH and are designed for key-based login (password-only SSH can block unattended runs).
\textbf{SSH bootstrap note:} run \texttt{scripts/rpi/setup\_ssh\_key.sh} to create a dedicated keypair under \texttt{scripts/rpi/keys/}. Pi helper scripts automatically use \texttt{scripts/rpi/keys/monad\_rpi5\_ed25519} when present (unless \texttt{SSH\_IDENTITY\_FILE} is explicitly set).

\subsection{Reference real-device smoke command}
\begin{verbatim}
SMOKE_PROFILE=rssi RSSI_DURATION_S=30 RSSI_INTERVAL_S_PI=2 \
RUN_PI=true INJECT_HYPOTHETICAL_METRICS=false DO_BUILD=false \
scripts/smoke/v3_end_to_end_smoke.sh
\end{verbatim}

\subsection{Reference wireless-spec command (recommended short real run)}
\begin{verbatim}
SMOKE_PROFILE=wireless_spec WIRELESS_SPEC_DURATION_S=30 \
WIRELESS_SPEC_INTERVAL_S_PI=3 WIRELESS_SPEC_BLE_TIMEOUT_S=12 \
WIRELESS_SPEC_CSI_TIMEOUT_S=15 WIRELESS_SPEC_ARTIFACT_STRIDE=0 \
RUN_PI=true INJECT_HYPOTHETICAL_METRICS=false DO_BUILD=false \
scripts/smoke/v3_end_to_end_smoke.sh
\end{verbatim}

If a real CSI collector is available on Pi, append collector variables:
\begin{verbatim}
CSI_COLLECTOR_CMD='sudo feitcsi --frequency 5180 --channel-width 80 --format HESU \
  --output-file /tmp/csi.dat -v' CSI_OUTPUT_PATH=/tmp/csi.dat
\end{verbatim}
\textbf{Important:} for Wi-Fi-only Pi control-plane, set \texttt{ALLOW\_WIFI\_DISRUPTIVE\_CSI=false} (default) unless you accept SSH drops during capture.

\subsection{Reference 2-minute full command}
\begin{verbatim}
SMOKE_PROFILE=full FULL_DURATION_S=120 FULL_INTERVAL_S_PI=2 \
FULL_ARTIFACT_STRIDE=0 REPORT_RPC_TIMEOUT_S=120 \
CONTROL_RPC_TIMEOUT_S=30 \
FULL_BLE_TIMEOUT_S=10 FULL_CSI_TIMEOUT_S=10 \
RUN_PI=true INJECT_HYPOTHETICAL_METRICS=false DO_BUILD=false \
scripts/smoke/v3_end_to_end_smoke.sh
\end{verbatim}

\section{Observability Queries}
\subsection{Grafana dashboard (provisioned)}
This repository provisions a Grafana dashboard focused on Raspberry Pi wireless smoke runs:
\begin{itemize}
  \item dashboard UID: \texttt{monad-fleet-wireless-v3}
  \item dashboard URL path: \texttt{/d/monad-fleet-wireless-v3/monad-fleet-wireless-sensing-pi-smoke-v3}
  \item provisioning files:
  \begin{itemize}
    \item \texttt{observability/grafana/provisioning/dashboards/monad-fleet.yml}
    \item \texttt{observability/grafana/provisioning/dashboards/monad-fleet-wireless-sensing-v3.json}
  \end{itemize}
\end{itemize}

\textbf{Panels include:} Wi-Fi RSSI, BLE advertisement count, Wi-Fi AP count, Wi-Fi connected status,
BLE scan status, CSI frame count, artifact upload counters, command health, full metric snapshot table,
and service counters.

\subsection{PromQL examples}
\begin{verbatim}
# Per-device totals
monad_fleet_metric_value{device_id="02:42:ac:14:00:04",metric="wifi_ap_total"}
monad_fleet_metric_value{device_id="02:42:ac:14:00:04",metric="ble_adv_total"}
monad_fleet_metric_value{device_id="02:42:ac:14:00:04",metric="csi_frames_total"}

# Real-device Wi-Fi indicators
monad_fleet_metric_value{device_id="2c:cf:67:80:f5:d9",metric="wifi_avg_rssi_dbm"}
monad_fleet_metric_value{device_id="2c:cf:67:80:f5:d9",metric="wifi_link_quality"}
monad_fleet_metric_value{device_id="2c:cf:67:80:f5:d9",metric="wifi_connected"}
monad_fleet_metric_value{device_id="2c:cf:67:80:f5:d9",metric="wifi_rssi_sample_index"}
monad_fleet_metric_value{device_id="2c:cf:67:80:f5:d9",metric="wifi_tx_bitrate_mbps"}
monad_fleet_metric_value{device_id="2c:cf:67:80:f5:d9",metric="wifi_rx_bitrate_mbps"}
monad_fleet_metric_value{device_id="2c:cf:67:80:f5:d9",metric="wifi_channel"}
monad_fleet_metric_value{device_id="2c:cf:67:80:f5:d9",metric="wifi_freq_mhz"}
monad_fleet_metric_value{device_id="2c:cf:67:80:f5:d9",metric="wifi_tx_retries"}
monad_fleet_metric_value{device_id="2c:cf:67:80:f5:d9",metric="wifi_if_rx_bytes"}
monad_fleet_metric_value{device_id="2c:cf:67:80:f5:d9",metric="wifi_if_tx_bytes"}

# Device status
monad_fleet_metric_value{device_id="2c:cf:67:80:f5:d9",metric="device_cpu_temp_c"}
monad_fleet_metric_value{device_id="2c:cf:67:80:f5:d9",metric="device_load1"}
monad_fleet_metric_value{device_id="2c:cf:67:80:f5:d9",metric="device_uptime_s"}
monad_fleet_metric_value{device_id="2c:cf:67:80:f5:d9",metric="device_mem_available_bytes"}

# Verify RSSI changed over interval window
changes(monad_fleet_metric_value{device_id="2c:cf:67:80:f5:d9",metric="wifi_avg_rssi_dbm"}[10m])

# BLE execution signal
monad_fleet_metric_value{device_id="2c:cf:67:80:f5:d9",metric="ble_scan_ok"}

# CSI collector status
monad_fleet_metric_value{device_id="2c:cf:67:80:f5:d9",metric="csi_collector_configured"}
monad_fleet_metric_value{device_id="2c:cf:67:80:f5:d9",metric="csi_capture_ok"}
monad_fleet_metric_value{device_id="2c:cf:67:80:f5:d9",metric="csi_capture_disabled"}
monad_fleet_metric_value{device_id="2c:cf:67:80:f5:d9",metric="csi_frames_total"}
monad_fleet_metric_value{device_id="2c:cf:67:80:f5:d9",metric="csi_output_files_count"}
monad_fleet_metric_value{device_id="2c:cf:67:80:f5:d9",metric="csi_output_bytes_total"}

# Low-resource sender sample
monad_fleet_metric_value{device_id="lowlevel-board-01"}

# Service counters
monad_fleet_reports_total
monad_fleet_events_total
monad_fleet_artifact_uploads_total
\end{verbatim}

\subsection{Interpretation notes}
\begin{itemize}
  \item \texttt{ble\_scan\_ok=1} means BLE scan command executed successfully; device discovery count may still be zero.
  \item \texttt{wifi\_avg\_rssi\_dbm} and \texttt{wifi\_link\_quality} are best-effort metrics based on available system tools.
  \item \texttt{wifi\_tx/rx\_bitrate\_mbps}, \texttt{wifi\_channel}, and interface counters provide richer Wi-Fi sensing context for RF/environment interpretation.
  \item \texttt{device\_*} metrics provide system-state context (thermal/load/memory/uptime) for cross-checking measurement stability.
  \item \texttt{wifi\_rssi\_sample\_index} tracks the latest sample ordinal within the configured sampling run.
  \item \texttt{csi\_frames\_total} may be zero if CSI collector is not configured or does not expose frame counts; check \texttt{csi\_collector\_configured} and \texttt{csi\_capture\_ok}.
  \item \texttt{csi\_capture\_disabled=1} indicates CSI execution was intentionally disabled by configuration (safety mode), not a collector failure.
\end{itemize}

\section{Compatibility and Versioning}
\begin{itemize}
  \item preferred interface for new devices: \texttt{fleet.v2}
  \item \texttt{fleet.v1} remains served for backward compatibility
  \item unknown fields/enums should be treated as forward-compatible additions
\end{itemize}

\section{Implementation Profile Decisions}
\begin{itemize}
  \item \texttt{PublishEvents} is optional and disabled by default in this deployment.
  \item if report events are empty, server may synthesize summary run event.
  \item if \texttt{Event.event\_id} is missing, server may synthesize an id.
  \item metric forwarding target can change at Prometheus \texttt{remote\_write} without changing the device gRPC contract.
  \item this deployment sets \texttt{METRICS\_EXCLUDE\_PREFIXES=csi\_} on Fleet service, so CSI metrics are not exported to Grafana.
  \item artifact upload to eLabFTW is implemented through \texttt{/ingest/v1/artifacts}; upload can be disabled per-agent if needed.
  \item sampled Wi-Fi runs throttle artifact fan-out by default (first+last sample artifact retention) to prevent oversized report processing on long sampling windows.
  \item agent-side report replay timeout is configurable via \texttt{REPORT\_RPC\_TIMEOUT\_S} and defaults to a higher value for Pi smoke runs.
  \item control-plane RPC timeouts can be tuned globally via \texttt{CONTROL\_RPC\_TIMEOUT\_S} or per method via \texttt{HELLO\_RPC\_TIMEOUT\_S}, \texttt{ASSIGNMENT\_RPC\_TIMEOUT\_S}, \texttt{POLICY\_RPC\_TIMEOUT\_S}, and \texttt{ACK\_PREPARED\_RPC\_TIMEOUT\_S}.
  \item CSI acquisition uses an external collector hook in this profile; collector integration is device/driver specific.
  \item Pi smoke runner applies a safety gate for Wi-Fi-only control-plane and can force \texttt{DISABLE\_CSI\_CAPTURE=true} unless \texttt{ALLOW\_WIFI\_DISRUPTIVE\_CSI=true}.
  \item very large CSI binaries should use object storage references when they exceed configured artifact limits.
\end{itemize}

\section{Deployment Defaults (Current)}
\begin{itemize}
  \item gRPC listen: \texttt{0.0.0.0:50060}
  \item HTTP listen: \texttt{0.0.0.0:9108}
  \item eLabFTW API base: \texttt{https://web/api/v2}
  \item fleet tag: \texttt{fleet}
  \item policy keys: \texttt{fleet.policy,policy,fleet\_policy,policy\_json}
  \item live events: disabled
  \item artifact upload endpoint max bytes: \texttt{ARTIFACT\_MAX\_BYTES=20971520}
  \item default spool root: \texttt{./data}
  \item sent retention: 14 days
  \item Prometheus: \texttt{http://localhost:9090}
  \item Mimir: \texttt{http://localhost:9009}
  \item Grafana: \texttt{http://localhost:3000} (\texttt{admin/admin} unless changed)
\end{itemize}

\section{Integrator Checklist}
Before declaring a new device implementation production-ready, confirm:
\begin{enumerate}
  \item agent sends stable \texttt{agent\_id} and \texttt{Hello} metadata
  \item agent follows prepare/measure/report lifecycle with retries
  \item idempotency keys are implemented (\texttt{preparation\_id}, \texttt{run\_id}, \texttt{event\_id})
  \item policy typed commands (\texttt{WIFI\_SCAN}/\texttt{BLE\_SCAN}/\texttt{CAPTURE\_CSI}) are handled or safely reported as unsupported
  \item run artifacts are persisted locally and published as \texttt{ArtifactRef.uri}, and optionally uploaded to eLabFTW via \texttt{/ingest/v1/artifacts}
  \item numeric metrics are exported via \texttt{summary\_metrics} and/or \texttt{event.metrics}
  \item smoke test validates report acceptance, metrics visibility, and eLabFTW artifact uploads
\end{enumerate}

\end{document}
